\section{The sixteen one-bit update rules}\label{sec:wirings}

The one-bit $\cxor$ model fixes the address calculation and residual emitter, leaving only the
cell update to vary.  If the update can inspect only the old state $s$ and current source bit
$x$, it is a function $w:\F^2\to\F$.  Four input pairs must each be assigned one output bit, so
there are
\[
  2^{|\F^2|}=2^4=16
\]
rules.  This is the complete update space under the one-cell, one-bit $\cxor$ restriction.

We name a rule by reading
\[
  w(0,0)\,w(0,1)\,w(1,0)\,w(1,1)
\]
as a four-bit binary integer, most significant bit first.  Thus the projection $w(s,x)=x$ has
truth table $0101$ and is $\W5$, while exclusive-or $w(s,x)=s+x$ has table $0110$ and is
$\W6$.

Every two-variable Boolean function has a unique algebraic normal form (ANF)
\begin{equation}\label{eq:anf-new}
  w(s,x)=d_0+d_1s+d_2x+d_3sx,
  \qquad d_i\in\F .
\end{equation}
Indeed, evaluating at $(0,0),(1,0),(0,1),(1,1)$ successively determines
\begin{equation}\label{eq:anf-coeff}
  d_0=w(0,0),\quad d_1=w(1,0)+d_0,\quad d_2=w(0,1)+d_0,
  \quad d_3=w(1,1)+d_0+d_1+d_2.
\end{equation}
This standard representation \cite{carlet2010boolean} exposes the classification needed below:
$d_3=0$ gives an \emph{affine} update, while $d_3=1$ couples the retained state to the observed
data.  In the language of Post's classification, the eight rules with $d_3=0$ constitute the
affine class \cite{post1941}; no further lattice machinery is needed here.

\begin{table}[ht]
\centering\small
\begin{tabular}{@{}clll@{\qquad}clll@{}}
\toprule
rule & truth & ANF & kind & rule & truth & ANF & kind\\
\midrule
$\W0$  & \bstr{0000} & $0$                 & affine & $\W8$  & \bstr{1000} & $1+s+x+sx$ & nonlinear\\
$\W1$  & \bstr{0001} & $sx$                & nonlinear & $\W9$  & \bstr{1001} & $1+s+x$ & affine\\
$\W2$  & \bstr{0010} & $s+sx$              & nonlinear & $\W{10}$ & \bstr{1010} & $1+x$ & affine\\
$\W3$  & \bstr{0011} & $s$                 & affine & $\W{11}$ & \bstr{1011} & $1+x+sx$ & nonlinear\\
$\W4$  & \bstr{0100} & $x+sx$              & nonlinear & $\W{12}$ & \bstr{1100} & $1+s$ & affine\\
$\W5$  & \bstr{0101} & $x$                 & affine & $\W{13}$ & \bstr{1101} & $1+s+sx$ & nonlinear\\
$\W6$  & \bstr{0110} & $s+x$               & affine & $\W{14}$ & \bstr{1110} & $1+sx$ & nonlinear\\
$\W7$  & \bstr{0111} & $s+x+sx$            & nonlinear & $\W{15}$ & \bstr{1111} & $1$ & affine\\
\bottomrule
\end{tabular}
\caption{The complete update-rule space.  Addition and multiplication are in $\F$.  ``Nonlinear''
means $d_3=1$ in ANF; the resulting recurrence has data-dependent state retention.}
\label{tab:new-wirings}
\end{table}

Two rules deserve names beyond their indices:
\begin{align}
  \W5(s,x)&=x,
  &\text{last-observation update,}\label{eq:w5-rule}\\
  \W6(s,x)&=s+x,
  &\text{running-parity update.}\label{eq:w6-rule}
\end{align}
$\W5$ is the rule used by the earlier $\cxor$ implementation.  It discards the old cell after
every observation and stores the successor just seen.  $\W6$ instead writes $s+x$, which is also the emitted residual
$r$ at that step.  Along repeated visits to one context it accumulates the parity of all observed
successors.  Although the same emitter permits reading its state as a hard prediction, its update
does not learn ``what followed this context last'' and is not naturally a next-symbol predictor.
It enters because it is the algebraic counterpart of $\W5$, not because it shares $\W5$'s modelling
interpretation.

All sixteen rules use the same causal emitter and are therefore invertible by
Theorem~\ref{thm:causal-inverse}.  The classification below asks what the update changes, rather
than using invertibility to choose an update.
